% HW3 brief: COMS E6998-019, Design of Production Agentic Systems.
% Build: cd briefs && tectonic -X compile -Z shell-escape -Z shell-escape-cwd=. hw3.tex
% minted needs pygmentize on PATH (e.g. uv tool install pygments).
\documentclass[9pt,letterpaper]{article}
\usepackage[extreme]{savetrees}
\usepackage[skip=\baselineskip,indent=0pt]{parskip}
\usepackage{booktabs,tabularx,array,enumitem,xcolor}
\usepackage[cache=false]{minted}
\usepackage[breakable]{tcolorbox}
\linespread{1.15}
\tcbuselibrary{skins} % the preamble's `enhanced' key needs this library
\usepackage{xltabular} % rubric tables may break across pages
\usepackage[hidelinks]{hyperref}
\setlist{nosep,leftmargin=*}
\renewcommand\tabularxcolumn[1]{>{\raggedright\arraybackslash}p{#1}}
\setminted{fontsize=\small,breaklines,frame=single,framerule=0.3pt,framesep=2mm}
\usepackage{needspace}
% Listings never split across pages. Each listing and its title sit in a box.
% The box moves to the next page whole if it does not fit.
\usepackage{etoolbox}
\BeforeBeginEnvironment{minted}{\par\smallskip\noindent\begin{minipage}{\linewidth}}
\AfterEndEnvironment{minted}{\end{minipage}\par\smallskip}
\newcommand\filelisting[3]{% title, lexer, file
  \par\smallskip\noindent\begin{minipage}{\linewidth}%
  \centerline{\small #1}\smallskip\inputminted{#2}{#3}\end{minipage}\par\smallskip}
\newtcolorbox{mustbox}[1]{enhanced,breakable,colback=red!3,colframe=red!60!black,
  title={#1},fonttitle=\bfseries}
% Long listings are floats. One float can take most of a page and share it with text.
\renewcommand\topfraction{0.9}
\renewcommand\textfraction{0.07}
\renewcommand\floatpagefraction{0.85}
\title{HW3: Deploying and Recovering an Agent on Kubernetes}
\author{COMS E6998 Section 019, Fall 2026 · Design of Production Agentic Systems\\
  Instructor: Rahul Krishna · TA: In Keun Kim}
\date{Released Friday, November 20, 2026 · Due Monday, December 14, 2026, before 11:59 p.m.}
\begin{document}
% savetrees sets \looseness=-1 on every paragraph.
% That setting breaks ragged table cells early, so keep normal line breaks.
\makeatletter\markeverypar{\the\everypar}\makeatother
\maketitle

% Inline code: prints its argument verbatim in typewriter type (no # or % inside).
\newcommand\code[1]{\texttt{\detokenize{#1}}}
% File paths: typewriter type, breakable after / . _ (no spaces inside).
\newcommand\file{\nolinkurl}

\section{Overview}\label{sec:overview}

HW3 asks two questions:
\begin{enumerate}
\item What does it take to deploy your agent on a Kubernetes cluster, and to
  resume it after a failure with correct checkpoints?
\item What does your Kubernetes deployment guarantee, and what did your HW2
  OpenShell sandbox guarantee that Kubernetes does not?
\end{enumerate}

\begin{itemize}
\item \textbf{Deploy on KIND.} Deploy your HW2 agent, its MCP server, and
  Postgres on a KIND cluster.
  \begin{itemize}
  \item You write all the manifests (Section~\ref{sec:build}).
  \item Staff deploy your stack from your README on the staff cluster
    (Sections~\ref{sec:readme} and~\ref{sec:cluster}).
  \end{itemize}

\item \textbf{Durable state.} Checkpoint the run to Postgres with LangGraph
  \code{PostgresSaver}.
  \begin{itemize}
  \item Resume the same run after a kill (Section~\ref{sec:state}).
  \item Staff kill your deployment during full scans (Section~\ref{sec:tested}).
  \end{itemize}

\item \textbf{Kubernetes guardrails.} Give the agent Pod a
  \code{securityContext}, writable volumes only where the agent needs them, and
  a NetworkPolicy.
  \begin{itemize}
  \item Section~\ref{sec:build} gives the objects and the settings.
  \end{itemize}

\item \textbf{Two test commands.} Add \code{--test-get-file} and
  \code{--test-access-host} to your agent.
  \begin{itemize}
  \item Each command tries one denied action. The expected result is a loud
    failure (Section~\ref{sec:testcmds}).
  \end{itemize}

\item \textbf{The report.} Explain how the agent recovers from each kill.
  Compare the Kubernetes guardrails with your HW2 OpenShell sandbox.
  \begin{itemize}
  \item Section~\ref{sec:report} gives the two parts of the report.
  \end{itemize}
\end{itemize}

HW3 keeps the HW2 remediation. The agent patches the alerts that it labels 1
and checks each patch. You design and implement the system behind these
interfaces.

\begin{center}
\begin{tabularx}{\linewidth}{@{}lX@{}}
\toprule
Released & Friday, November 20, 2026, after class \\
Due & Monday, December 14, 2026, before 11:59 p.m. \\
Weight & 20\% of the course grade, marked out of 20 (Section~\ref{sec:rubric}) \\
Format & Team assignment. One submission per team. \\
Late days & 6 per student per semester, at most 3 on one assignment. Each late day extends the deadline by 24 hours. A late HW3 can arrive by Thursday, December 17, 2026. \\
\bottomrule
\end{tabularx}
\end{center}

\textbf{The model.} HW3 uses the university OpenAI-compatible endpoint. Set
\code{model.provider} to \code{"openai-compatible"}. You know the base URL and
the model name, so this brief and the examples use placeholders. Staff give each
team its own key, with a \$25 cap.

The cap on the key is the hard limit.
\code{budget.max_tokens} in the config is a soft limit, with its HW2 rules.

\textbf{The teaching dates.} The lectures teach the Kubernetes topics of
HW3 on three dates:
\begin{itemize}
\item \textbf{November 20.} Controllers, Pods, init containers, and Pod
  deletes. The Pod life cycle, probes, and termination. \code{emptyDir}
  compared with a PVC. StatefulSets and headless Services. A demo of
  \code{PostgresSaver} after a Postgres Pod delete.
\item \textbf{December 4.} Environment variables, ConfigMaps, and Secrets. Jobs,
  \code{restartPolicy}, and \code{backoffLimit}. The Pod \code{securityContext}.
\item \textbf{December 11.} CNI plugins, NetworkPolicy, Calico, and an
  integration studio.
\end{itemize}
Section~\ref{sec:devpath} tells you how to start work before December 4.

\section{What you build}\label{sec:build}

\subsection{The HW2 baseline}\label{sec:baseline}

Start from your own HW2 MCP agent, MCP server, and remediation. You can repair
them. Your README names the commit of your HW2 baseline
(Section~\ref{sec:deliverables}). HW3 does not use the code-action agent or the
OpenShell sandbox. We do not regrade HW2.

\subsection{The objects}\label{sec:objects}

You write each of these objects in full:

\begin{center}
\small
\begin{tabularx}{\linewidth}{@{}>{\raggedright\arraybackslash}p{0.47\linewidth}X@{}}
\toprule
Object & What it teaches \\
\midrule
Postgres StatefulSet, headless Service, PVC, Secret & Storage, stable identity, Secrets \\
MCP server Deployment and Service, with readiness and liveness probes & Service addressing, probes \\
ConfigMap for \code{config.toml}, and an input PVC for the targets & Configuration, volumes \\
Secret for the team key, passed as an environment variable & Secrets \\
Agent Job, with an init container that waits for Postgres & Jobs, restart policy, \code{backoffLimit} \\
\code{securityContext} on the agent Pod & Non-root, no capabilities, read-only root file system \\
Writable volumes for the agent & Where patches and scratch files go \\
NetworkPolicy for the agent Pod & Egress only to MCP, Postgres, DNS, and the university endpoint \\
\bottomrule
\end{tabularx}
\end{center}

\textbf{MCP transport.} Your HW2 MCP server uses stdio. In HW3, the server is a
separate Deployment, so the agent connects to it over the network. Change the
server to the MCP streamable HTTP transport. In the MCP Python SDK, call
\code{MCPServer.run(transport="streamable-http")} with \code{host} and
\code{port}. The client uses \code{streamable_http_client}.

The default
\code{host} is \code{127.0.0.1}, and other Pods cannot connect to it.
\code{[mcp] url} gives the agent the server URL (Section~\ref{sec:config}).

\subsection{Volumes}\label{sec:volumes}

A read-only root file system does not stop the agent from writing patches. The
image stays read-only, and the agent writes only to mounted volumes:

\begin{center}
\small
\begin{tabularx}{\linewidth}{@{}llX@{}}
\toprule
Mount & Volume & Content \\
\midrule
\code{/work} & PVC & The disposable copy of the target, the patches, and the trace. A restart keeps them. \\
\code{/tmp} & \code{emptyDir} & Scratch files. A Pod delete removes them. \\
\code{/inputs} & Read-only PVC & The pinned target source. The agent copies the allowed files to \code{/work}. \\
\bottomrule
\end{tabularx}
\end{center}

Your report states which files survive a process kill and which files survive a
Pod delete (Section~\ref{sec:report}).

\subsection{Settings}\label{sec:settings}

\begin{itemize}
\item \textbf{The agent is a Job, not a Deployment.} A batch run has an end. A
  Deployment starts a new Pod each time the agent exits.
\item \textbf{\code{restartPolicy: OnFailure} and a \code{backoffLimit} of 10
  or more.} A process kill then restarts the container in the same Pod. The Job
  stops only after many failures.
\item \textbf{An init container runs \code{pg_isready} until Postgres is ready.}
  A readiness probe on Postgres does not delay the agent start. The init
  container does.
\item \textbf{Postgres is a separate StatefulSet.} A kill of the agent Pod must
  not stop the database. The database keeps its PVC across a Pod delete.
\end{itemize}

\subsection{Network policy}\label{sec:netpol}

The API server stores a NetworkPolicy, and the CNI plugin enforces it. A CNI with
no policy support ignores the policy and gives no warning. Without enforcement,
\code{--test-access-host} cannot fail.

% OPEN: feasibility task 3. The CNI choice has no result on the build day.
\textbf{CNI.} kindnet or Calico. Staff publish the choice by December 4.
Section~\ref{sec:cni} gives a deny-all check that you can run.

The egress rule of the agent Pod allows only four destinations: the MCP server,
Postgres, DNS, and the university endpoint. A NetworkPolicy has no hostname
field. It selects Pods by label, and other hosts by IP block and port.

\begin{mustbox}{Deployment rules}
\begin{itemize}
\item Your team writes all the manifests.
\item The manifests are plain Kubernetes YAML, in one file or with Kustomize.
  Staff do not accept Helm or other template tools.
\item No committed file holds a Secret value.
\item The agent runs as a Job.
\item The root file system of the agent container is read-only.
\end{itemize}
\end{mustbox}

\section{Targets}\label{sec:targets}

Scan each of the three HW1 targets in full, at the HW1 pins. Use Bandit 1.9.4 at
MEDIUM or higher severity. HW1 and HW2 used 49 alerts in total. HW3 uses the full
scans. The staff \code{[scan]} blocks are in \file{briefs/scan/}. Each alert count
below is the count after the cap.

\begin{center}
\small
\newcommand\repo[2]{\begin{tabular}[t]{@{}l@{}}\footnotesize\url{#1}\\\ttfamily\footnotesize #2\end{tabular}}
\begin{tabularx}{\linewidth}{@{}lXlr@{}}
\toprule
Target & Repository and commit & \code{[scan]} block & Alerts \\
\midrule
OWASP BenchmarkPython & \repo{https://github.com/OWASP-Benchmark/BenchmarkPython}{f1291485808b66e20ddb6b01b10dc71b3df8c8ba} & \file{owasp-benchmark-python.toml} & 297 \\
Radicale & \repo{https://github.com/Kozea/Radicale}{eff8027f3dc4910be1659ee71f4b4a454ade5a8c} & \file{radicale.toml} & 10 \\
SGLang & \repo{https://github.com/sgl-project/sglang}{bbe9c7eeb520b0a67e92d133dfc137a3688dc7f2} & \file{sglang.toml} & 21 \\
\midrule
Total & & & 328 \\
\bottomrule
\end{tabularx}
\end{center}

\textbf{The cap.} Each block sets \code{max_alerts = 500}. If a scan gives more
than 500 alerts, the agent keeps the first 500. The order is HIGH before MEDIUM,
then the HW2 row order: \code{path}, \code{line} as an integer, and
\code{test_id}.

A target with 500 alerts or fewer keeps all its alerts. At these
pins, no target reaches the cap. OWASP has the largest count, 297. The cap stays
the rule.

\filelisting{\file{briefs/scan/owasp-benchmark-python.toml}, the OWASP
\code{[scan]} block}{toml}{scan/owasp-benchmark-python.toml}

\textbf{The answer key.} OWASP's answer key, \code{expectedresults-0.1.csv}, is
open to you and never to your agent, as in HW1 and HW2. In HW3, the answer key
is not on \code{/inputs}.

\textbf{Remediation.} The HW2 remediation rules are unchanged. The agent patches
each alert that it labels 1. It accepts a patch only if the patch passes the
three checks: \code{applies}, \code{touches_flagged_line}, and
\code{clean_rescan}. Each row keeps the HW2 \code{remediation} field. The HW2
brief gives the details.

\textbf{Cost.} A full run on all three targets must fit in the \$25 key.
% OPEN: feasibility task 2. The per-alert token cost has no result on the build day.
The per-alert token cost on the university endpoint is to follow. Measure your
own cost on a small run first. Then set \code{budget.max_tokens} so that a full
run stays under the cap.

\section{Durable state and resume}\label{sec:state}

A kill must not start a new run. After each kill, the agent loads its last
checkpoint and continues the same run. This section gives the three stores that
make this possible: the checkpoints, the results table, and the ledger. All three
are in Postgres, on the Postgres PVC.

\subsection{Checkpoints}\label{sec:checkpoints}

LangGraph \code{PostgresSaver} stores the checkpoints. The agent calls
\code{setup()} at each start. The call makes the checkpoint tables if they do not
exist. The \code{thread_id} of the run is \code{[run] id} from the config. Set
\code{durability="sync"}, so that each checkpoint write completes before the
next step starts. Your README states the setting and the reason.

\begin{minted}{python}
from langgraph.checkpoint.postgres import PostgresSaver

with PostgresSaver.from_conn_string(conn_string) as saver:   # Postgres from [postgres]
    saver.setup()                                             # make the tables once
    app = graph.compile(checkpointer=saver)
    run_id = config["run"]["id"]                              # the same ID after each kill
    thread = {"configurable": {"thread_id": run_id}}
    first = saver.get_tuple(thread) is None                   # no checkpoint yet
    for update in app.stream(start_state if first else None, thread, durability="sync"):
        ...
\end{minted}

A new UUID at each start is a common error. Then a restart starts a new run, with
no checkpoint. The agent scans all the alerts again and charges the key again.

\subsection{The results table}\label{sec:results}

The agent writes each finished alert row to a results table in Postgres. The
primary key is \code{alert_id}, the HW1 alert ID. On a resume, the agent skips
each alert that has a row. No \code{alert_id} appears twice. The agent writes
\code{report.json} from this table at the end of the run.

\subsection{The ledger}\label{sec:ledger}

The ledger is a Postgres table that records each model call in two phases:
\begin{enumerate}
\item Before the call, the agent writes a \code{call_started} row.
\item After the reply, the agent writes a \code{call_finished} row with the usage
  that the provider reported.
\end{enumerate}
An open row is a \code{call_started} row with no \code{call_finished} row for
the same \code{call_id}. A kill during a model call leaves an open row. The
provider can charge the key for that call, but the agent has no record of the
usage. Your report counts the open rows as possible charges with no record.
Section~\ref{sec:ledgerfmt} gives the export format.

\subsection{What a restart keeps}\label{sec:keeps}

\begin{center}
\small
\begin{tabularx}{\linewidth}{@{}lXX@{}}
\toprule
Item & After a process kill & After a Pod delete \\
\midrule
Checkpoint (Postgres) & kept & kept \\
Results table (Postgres) & kept & kept \\
Ledger (Postgres) & kept & kept \\
\code{/work} (PVC) & kept & kept. The new Pod mounts the same PVC. \\
\code{/tmp} (\code{emptyDir}) & kept. The container restarts in the same Pod. & lost \\
Trace (on \code{/work}) & kept. The agent appends to it. & kept. The agent appends to it. \\
\bottomrule
\end{tabularx}
\end{center}

A kill loses the memory of the agent process in both cases. Work after the last
checkpoint can repeat (Section~\ref{sec:trace}).

\begin{mustbox}{Resume rules}
\begin{itemize}
\item The run ID is the same after each kill.
\item No \code{alert_id} appears twice in the results table or in
  \code{report.json}.
\item The ledger never resets. It continues from its rows before the kill.
\item The agent writes a \code{resume} event at each start.
\end{itemize}
\end{mustbox}

\section{Interface, output, trace, and ledger}\label{sec:interface}

The evaluator uses your configuration to run the system without reading your
code. The \code{[scan]} block specifies the target files. \code{[run] id} names
the run. \code{[postgres]} and \code{[mcp]} give the two services, and
\code{[budget]} sets the limits to check. Grading cannot proceed if your
configuration fails \code{config.schema.json} validation.

\subsection{Command line}\label{sec:cli}

\begin{minted}{bash}
your-agent --input <dir> --output <dir> --config <file>
your-agent --config <file> --test-get-file <path>
your-agent --config <file> --test-access-host <host:port>
\end{minted}
The scan command is the HW2 command. It writes \code{report.json}
(Section~\ref{sec:output}) and \code{trace.jsonl} (Section~\ref{sec:trace}) to
the output directory, on the \code{/work} PVC. Inside the Pod, \code{--input} is
the copy of the allowed files on \code{/work}, not \code{/inputs}. The two test
commands are new (Section~\ref{sec:testcmds}). Choose your own executable name in
place of \code{your-agent} and document it in the README.

\subsection{Config}\label{sec:config}

Your agent reads the configuration file that \code{--config} names. It is a TOML
file. After parsing, it must validate against \code{config.schema.json}. These
are the keys:

\begin{center}
\small
\begin{tabularx}{\linewidth}{@{}lX@{}}
\toprule
Key & Rule \\
\midrule
\code{[scan]} & Exactly one of the three blocks in \file{briefs/scan/}, unchanged. Each block has \code{max_alerts = 500}. \\
\code{model.provider} & \code{"openai-compatible"}. \\
\code{model.base_url}, \code{model.name} & The university endpoint and the model name. \\
\code{model.api_key_var} & The name of the environment variable that holds the key. Never the key. \\
\code{budget.max_tokens} & Integer $\geq 0$: the token budget for the whole run, across all alerts. A soft limit. \\
\code{budget.output_allowance} & Integer $\geq 1$: the maximum output tokens for each model call, sent to the provider. \\
\code{budget.max_steps} & Integer $\geq 1$: a failsafe, separate from the token budget. \\
\code{[run] id} & A string: the checkpoint \code{thread_id}. Staff set it at grading. \\
\code{[checkpoint] durability} & \code{"sync"}. \\
\code{[postgres]} & \code{host}, \code{port}, \code{database}, \code{user}, and \code{password_var}. \code{password_var} names an environment variable. \\
\code{[mcp] url} & The URL of the MCP Service, for example \code{http://mcp:8000/mcp}. \\
\bottomrule
\end{tabularx}
\end{center}
HW3 removes the HW2 keys \code{[agent] kind} and \code{[sandbox]}. You may add
keys for your own settings. Read secrets from environment variables. Never put
them in the configuration file.

Staff set \code{[run] id} from the environment variable \code{HW3_RUN_ID} at
each run (Section~\ref{sec:readme}). The base URL and the model name below are
placeholders.

\filelisting{\file{examples/config.example.toml}, a valid config for OWASP
BenchmarkPython}{toml}{../examples/config.example.toml}

\subsection{Output}\label{sec:output}

Your \code{report.json} must validate against \code{report.schema.json}.
Figure~\ref{fig:report} shows a cut-down run with 3 rows. It sets
\code{alert_count} to 3 to match its rows. A full OWASP run has 297 rows. Its
rationales are placeholders.

Follow these rules:
\begin{itemize}
\item \code{run_id} is the same string as \code{[run] id} in the config.
\item There is exactly one row for each in-scope Bandit result. Rows are never
  dropped or merged. \code{len(alerts)} equals \code{scan.alert_count}, the count
  after the cap.
\item \code{alert_id} keeps its HW1 form and is unique across rows. The agent
  scans the read-only \code{/inputs} copy, so the IDs do not change during a run.
\item \code{path} and every evidence \code{path} are repository-relative POSIX
  paths with no leading \code{./}. Bandit prints \code{./} for explicitly listed
  files, so strip it. Rows are in order of \code{path} as a plain string, then
  \code{line} as an integer, then \code{test_id}.
\item \code{status} is one of \code{classified}, \code{budget_exhausted} or \code{error}.
\item \code{label} is \code{1} or \code{0} when \code{status} is \code{classified}.
  Otherwise it is \code{null}.
\item \code{budget.used_tokens} is exactly the budget charge, which covers the
  whole run as in HW1. It is \code{input + output} as the provider reported them,
  summed over every \code{model_call}. \code{cache_read} and \code{cache_write}
  are never charged.
\item \code{target.commit} is the pinned commit of Section~\ref{sec:targets}.
  Give the agent the commit through a key of your own in the config, for
  example.
\item Every row has a \code{remediation} object with \code{status},
  \code{diff}, and \code{checks}, as in HW2. The \code{diff} is a unified diff,
  or an empty string if the agent proposed no patch. The \code{checks} object
  holds \code{applies}, \code{touches_flagged_line}, and \code{clean_rescan}.
  Each check is \code{true}, \code{false}, or \code{null} if the check did not run.
\item The remediation \code{status} is \code{accepted}, \code{rejected}, or
  \code{skipped}. An accepted patch has a non-empty diff and all three checks
  \code{true}. A rejected patch failed a check. Use \code{skipped} for every
  alert labelled 0 or left unclassified.
\end{itemize}

\begin{figure}[t]
\inputminted[fontsize=\scriptsize]{json}{../examples/report.example.json}
\caption{\file{examples/report.example.json}, a cut-down run with 3 OWASP rows.}\label{fig:report}
\end{figure}

\Needspace{14\baselineskip}
\subsection{Execution trace}\label{sec:trace}

The trace records every step of a run in order. It records whether the model or
the runtime made each decision. Implement the trace in the format below. For
reference, see Cyberbird's \file{cyberbird/plan_and_validation/trace.py}
(\code{Trace}, \code{EventKind}, \code{TerminalStatus}, \code{Usage}) and
\file{cyberbird/plan_and_validation/trajectory.py}. The second file converts a
trace into a Mermaid sequence diagram:
\par\noindent\begin{minipage}{\linewidth}\raggedright
\url{https://github.com/comse6998-019/cyberbird/blob/cf7e96c8bd9e74470b99b7b754bb72392427cfc6/cyberbird/plan_and_validation/trace.py}\\
\url{https://github.com/comse6998-019/cyberbird/blob/cf7e96c8bd9e74470b99b7b754bb72392427cfc6/cyberbird/plan_and_validation/trajectory.py}
\end{minipage}\par

Write the trace as follows:
\begin{itemize}
\item \code{<out dir>/trace.jsonl} holds one JSON object per line. It is
  append-only, on the \code{/work} PVC. Flush it after every event, so that the
  trace from before a kill stays.
\item Every event has \code{step} (an integer from 1, up by 1),
  \code{ts} (Unix time, float), \code{run_id}, \code{alert_id} (the alert of the
  event, or \code{null} for run-level events) and \code{kind}.
\item \code{kind} is one of the values in the table below, or \code{discovery}
  from the box after the table. Events may include more fields.
\end{itemize}

\begin{center}
\small
\begin{tabularx}{\linewidth}{@{}lp{0.25\linewidth}X@{}}
\toprule
\code{kind} & Written by & Fields \\
\midrule
\code{resume} & the runtime, once at each start, first & \code{attempt} (1 at the
  first start), \code{checkpoint_id} (\code{null} at the first start), and
  \code{next_alert_id} (\code{null} if no alert remains) \\
\code{model_call} & the node that called the model & \code{role}, \code{model},
  \code{usage} with \code{input} and \code{output} (integers, as the provider
  reported them. Add \code{cache_read} and \code{cache_write} if the provider
  reports them. No summed total.) Add \code{call_id}, the ledger
  \code{call_id} of the call. \\
\code{tool_request} & the dispatcher & \code{tool}, \code{args} \\
\code{tool_result} & the dispatcher & \code{tool}, \code{ok} (boolean), and
  \code{error} when \code{ok} is false \\
\code{routing} & the edge function that chose the next node & \code{decision},
  \code{by} (\code{model} or \code{runtime}) \\
\code{state_change} & any node, for updates not in the rows above & \code{node}.
  When an alert row is final (\code{classified}, \code{error}, or
  \code{budget_exhausted}), also \code{alert_status} and \code{label}. \\
\code{terminal} & the runtime, once, last & \code{status} (one of your terminal
  statuses, as in HW1), \code{usage_total} (an object with the same counters as \code{usage}) \\
\bottomrule
\end{tabularx}
\end{center}

Every trace must satisfy these invariants:
\begin{enumerate}
\item The file parses line by line, and every event has the fields for its
  \code{kind}.
\item \code{step} starts at 1 and goes up by exactly 1. After a restart,
  \code{step} continues from the last line in the file.
\item There is exactly one \code{terminal} event, at the end of the run, and it
  is the last line. A killed attempt writes no \code{terminal} event.
\item Each counter in \code{usage_total} equals the sum of that counter over all
  \code{model_call} events. \code{budget.used_tokens} in \code{report.json}
  equals exactly the sum of \code{input + output} over those events.
  \code{cache_read} and \code{cache_write} are recorded but never charged.
\item Work after the last checkpoint can repeat, and the trace keeps both
  copies. Each row of \code{report.json} has at least one \code{state_change}
  event with its \code{alert_id}, \code{alert_status}, and \code{label}. Each such
  event must agree with the last one for that alert, and the last one must agree
  with the row.
\item Every \code{tool_request} is immediately followed by its \code{tool_result}
  (the dispatcher handles one request at a time), before the next
  \code{model_call}.
\item The first event of each attempt is \code{resume}. The first
  \code{resume} has \code{attempt} 1, and each later \code{resume} has the
  \code{attempt} before it plus 1.
\end{enumerate}

Each line must validate against \code{trace.schema.json} at the repository root.
Graders check the invariants across events separately, because a schema for one
event cannot express them.

\begin{mustbox}{HW3 trace and schema rules}
\begin{itemize}
\item Include \code{via} on each \code{tool_request} and \code{tool_result}. Its
  only value is \code{mcp}.
\item Include \code{refused_by} on each refused \code{tool_result}. Its value
  names the layer that refused the call: \code{dispatcher} or \code{server}.
\item After \code{tools/list}, the MCP host writes a \code{discovery} event with
  \code{server} (the server name) and \code{tools} (the tool names offered to
  the model). The host writes it again after each resume.
\item For remediation, write a \code{state_change} event with
  \code{node: remediation} and \code{checks}, which holds the three check
  results. Do not include \code{alert_status} or \code{label} in this event.
\item Use the HW3 versions of \code{config.schema.json},
  \code{report.schema.json}, \code{trace.schema.json}, and
  \code{ledger.schema.json}. The config needs \code{[run]}, \code{[checkpoint]},
  \code{[postgres]}, and \code{[mcp]}. The trace schema accepts \code{resume}.
  The report needs \code{run_id}.
\end{itemize}
\end{mustbox}

\file{examples/trace.example.jsonl} shows one restart. The agent dies during a
model call on the first alert, writes a \code{resume} event, and finishes the run.
It is the same run as \file{examples/report.example.json} and
\file{examples/ledger.example.jsonl}, and it meets all seven invariants.

\filelisting{\file{examples/trace.example.jsonl}}{json}{../examples/trace.example.jsonl}

\subsection{Ledger}\label{sec:ledgerfmt}

Export the ledger table as JSON Lines, one row per line, in the order of
\code{ts}. Each line must validate against
\file{ledger.schema.json}.

\begin{center}
\small
\begin{tabularx}{\linewidth}{@{}lX@{}}
\toprule
Field & Rule \\
\midrule
\code{call_id} & One ID for each model call. One \code{call_started} row and at most one \code{call_finished} row share it. \\
\code{run_id} & The \code{run_id} of \code{report.json}. \\
\code{alert_id} & The alert of the call. \\
\code{attempt} & The \code{attempt} of the \code{resume} event before the call. \\
\code{status} & \code{call_started} or \code{call_finished}. \\
\code{ts} & Unix time, float. It goes up from row to row. \\
\code{usage} & On a \code{call_finished} row only: \code{input}, \code{output}, and the cache counters if the provider reports them. \\
\bottomrule
\end{tabularx}
\end{center}

\textbf{Open rows.} An open row is a \code{call_started} row with no
\code{call_finished} row for the same \code{call_id}. An open row is never
charged, and it has no \code{model_call} event in the trace.

\textbf{Reconciliation.} The sum of \code{input + output} over the
\code{call_finished} rows equals \code{budget.used_tokens} in \code{report.json}.
Each \code{call_finished} row has one \code{model_call} event with the same
\code{call_id} and \code{usage}.

In the example, \code{call-0002} is open. The kill came during that call.

\filelisting{\file{examples/ledger.example.jsonl}}{json}{../examples/ledger.example.jsonl}

\section{Test commands}\label{sec:testcmds}

The agent command line adds two commands. Each command tries one denied action,
prints the error, and exits with a nonzero code.

\begin{center}
\small
\begin{tabularx}{\linewidth}{@{}>{\raggedright\arraybackslash}p{0.2\linewidth}XX@{}}
\toprule
Command & Action & Expected loud failure in Kubernetes \\
\midrule
\code{--test-get-file <path>} & Read a path outside the allowed set, for example the OWASP answer key & \code{ENOENT} (not mounted) or \code{EACCES} (read-only mount, UID) \\
\code{--test-access-host <host:port>} & Open a connection to a host that the NetworkPolicy denies & Timeout or connection refused from the CNI \\
\bottomrule
\end{tabularx}
\end{center}

\textbf{Output.} On a denial, print one line to standard error,
\code{denied: <ERRNO or TIMEOUT> <detail>}, and exit with code 1. If the action
succeeds, print \code{allowed: <detail>} and exit with code 0. That result is a
guardrail failure.

\begin{minted}{text}
$ your-agent --config /config/config.toml --test-get-file /inputs/expectedresults-0.1.csv
denied: ENOENT /inputs/expectedresults-0.1.csv
$ your-agent --config /config/config.toml --test-access-host 203.0.113.7:443
allowed: connected to 203.0.113.7:443
\end{minted}

\textbf{Rules.}
\begin{itemize}
\item \code{--test-access-host} uses a 5 second connect timeout. A command that
  hangs with no output is not a loud failure.
\item The commands make no model call, no Postgres write, and no MCP call. They
  cost nothing.
\item Exit code 1 means denied. Exit code 0 means allowed.
\end{itemize}

\textbf{The test Pod.} A finished Job has no container for \code{kubectl exec},
so staff run the commands in a test Pod. Supply \file{k8s/agent-test-pod.yaml}.
The test Pod has the same image, labels, \code{securityContext}, and volumes as
the agent Job Pod. The labels select the same NetworkPolicy. Staff set the
arguments.

Do not run the commands in the HW2 sandbox again. Part 2 of the report uses the
HW2 evidence that you already have (Section~\ref{sec:report}).

\section{How staff test the deployment}\label{sec:tested}

You do not need to write a test harness. Staff kill your deployment during full
scans and run the two test commands in the cluster. This section gives the kill
points, the methods, and the pass criteria.

\subsection{Kill points}\label{sec:killpoints}

\begin{itemize}
\item Between two alerts.
\item During a tool call.
\item During a model call.
\end{itemize}
Staff can kill the agent at any time, including during a model call.

\subsection{Kill methods}\label{sec:killmethods}

\begin{center}
\small
\begin{tabularx}{\linewidth}{@{}>{\raggedright\arraybackslash}p{0.27\linewidth}>{\raggedright\arraybackslash}p{0.27\linewidth}X@{}}
\toprule
Method and command & Effect & What it tests \\
\midrule
Process kill: \code{kubectl exec <agent pod> -- kill -9 <agent pid>} & Kubernetes restarts the container in the same Pod & \code{restartPolicy: OnFailure}, checkpoint reload \\
Agent Pod delete: \code{kubectl delete pod <agent pod>} & The Job makes a new Pod & \code{emptyDir} loss, PVC reattach, \code{backoffLimit} \\
Postgres Pod delete: \code{kubectl delete pod <postgres pod>} & The StatefulSet makes a new Postgres Pod on the same PVC & The agent waits or retries, and loses no state \\
MCP Pod delete: \code{kubectl delete pod <mcp pod>} & The Deployment makes a new MCP Pod & The agent retries the tool call \\
\bottomrule
\end{tabularx}
\end{center}

% OPEN: feasibility task 5. The PID of the agent for kill -9 has no result on the build day.
The process kill names \code{<agent pid>}. The kernel ignores a \code{SIGKILL}
that a process in the container sends to PID 1. So \code{kill -9 1} kills the
agent only if a test shows it. Staff publish the exact command by December 4.

\subsection{Schedule}\label{sec:schedule}

Each run has four kills, one for each method. A seed selects the time of each
kill. Each of the three kill points occurs at least once in a run. Staff choose
the seed at grading time.

\subsection{Pass criteria}\label{sec:pass}

A run passes if all five criteria hold:
\begin{enumerate}
\item The run ID is the same before and after each kill.
\item \code{len(alerts)} equals the published alert count.
\item No \code{alert_id} appears twice in the results table.
\item The ledger continues from its value before the kill. It does not reset.
\item The trace shows the restart, the checkpoint that loaded, and the first
  alert after the resume.
\end{enumerate}

\subsection{The test commands}\label{sec:stafftests}

Staff run both commands in your test Pod, with paths and hosts that staff choose.
Each command must print a \code{denied:} line and exit with code 1.

\subsection{Your own kill runs}\label{sec:ownkills}

You can run the same kills on your own stack before you submit. Use the commands
in the table above. The brief gives no script.

\section{Report}\label{sec:report}

The report has no page limit. It has two parts.

\textbf{Part 1: deployment and recovery.}
\begin{enumerate}
\item For each kill point and each kill method, explain how your design
  handles the kill:
  \begin{itemize}
  \item which checkpoint loads
  \item which work the agent does again
  \item how the ledger treats an interrupted model call
  \item which files survive
  \end{itemize}
\item Explain the checkpoint design: what the checkpoint holds, the
  \code{durability} setting, and the ledger rules.
\item Name the fault windows that the staff tests do not cover.
\end{enumerate}

\textbf{Part 2: Kubernetes compared with the HW2 OpenShell sandbox.}
\begin{enumerate}
\item Explain how file denial and host denial work in the Kubernetes deployment
  and in the HW2 sandbox.
\item State what the HW2 OpenShell sandbox guaranteed that the Kubernetes
  deployment does not.
\item State what the Kubernetes deployment gives that the HW2 sandbox did not.
\end{enumerate}
For Part 2, cite the output of your own test commands in the cluster and your
own HW2 evidence.

\section{Rubric}\label{sec:rubric}

The assignment is marked out of 20, in four criteria of 5 points. The marks
assess the mechanisms and the evidence for them. They do not assess the
correctness of the labels.

{\small
\begin{xltabular}{\linewidth}{@{}>{\raggedright\arraybackslash}p{0.24\linewidth}r X@{}}
\toprule
Criterion & Points & What earns credit \\
\midrule
Systems design and claim & 5 & \textbf{3:} the Job, Postgres, \code{PostgresSaver},
  the ledger, and the \code{alert_id} key. \textbf{2:} the \code{securityContext},
  the writable volumes, and the NetworkPolicy. \\
Experimental design & 5 & \textbf{3:} the staff kill tests pass the five criteria
  of Section~\ref{sec:pass}. \textbf{2:} both test commands fail loudly in the
  cluster. \\
Evidence and reproducibility & 5 & \textbf{3:} staff deploy the app on the staff
  KIND cluster from the README alone. \textbf{2:} the traces, the ledger, and the
  results reconcile. \\
Interpretation and limitations & 5 & \textbf{2:} Part 1 explains each kill.
  \textbf{2:} the Part 2 comparison is correct. \textbf{1:} the report names the
  untested fault windows. \\
\midrule
Total & 20 & \\
\bottomrule
\end{xltabular}
}

Grading notes:
\begin{itemize}
\item A sound experiment can earn full credit without a confirmed hypothesis.
\item A missing mechanism is a defect. A description of why it is absent does
  not earn its points.
\item One problem gets one penalty. We do not mark the same problem twice.
\item Accuracy, precision, and recall do not affect grades.
\end{itemize}

\textbf{Staff procedure.} Staff make a KIND cluster from
\file{cluster/kind-config.yaml} and deploy your stack from your README alone.
Then staff run the kill schedule of Section~\ref{sec:schedule} on full scans, and
run the two test commands in your test Pod.

\section{README}\label{sec:readme}

Staff use your \code{README.md} to deploy your stack from a fresh clone, with no
other help. Staff deploy on a staff laptop and use only the files in your
repository. Document the following:
\begin{enumerate}
\item \textbf{Setup.} The tools and their versions: Docker, KIND, \code{kubectl},
  Python 3.12, Bandit 1.9.4, CLDK, the MCP SDK, LangGraph, and your other
  dependencies. List the environment variables that your stack reads, with
  names only and never values.
\item \textbf{The staff steps.} The exact commands, in order, from a fresh clone
  to a finished run. The steps are below.
\item \textbf{Durability.} The \code{durability} setting and the reason for it
  (Section~\ref{sec:checkpoints}).
\item \textbf{Configs.} One config for each target. The three configs are
  identical except for \code{[scan]}.
\item \textbf{The rubric index.} A table from each rubric criterion to the files
  and the trace lines that give its evidence.
\end{enumerate}

\subsection{The staff steps}\label{sec:staffsteps}

Your README tells staff how to do these seven steps:
\begin{enumerate}
\item Build each image with \code{docker build}. Supply one Dockerfile for each
  image that you build, for example the agent and the MCP server.
\item Load each image into the cluster with \code{kind load docker-image}. The
  staff cluster cannot pull from a team registry.
\item Make the two Secrets, for the team key and the Postgres password. No
  committed file holds a key.
\item Apply the manifests.
\item Copy the pinned target source to the input PVC.
\item Start a run. Give one command that reads the run ID from the environment
  variable \code{HW3_RUN_ID} and applies the agent Job.
\item Get \code{report.json}, \code{trace.jsonl}, and the ledger export.
\end{enumerate}

\textbf{Manifests.} Put plain Kubernetes YAML in a \file{k8s/} directory at the
root of your repository. Use one of two forms:

\begin{center}
\small
\begin{tabularx}{\linewidth}{@{}lXl@{}}
\toprule
Form & Files & Command that staff run \\
\midrule
One file & \file{k8s/app.yaml}, with all objects and \code{---} between them & \code{kubectl apply -f k8s/app.yaml} \\
Kustomize & Manifest files and a \file{k8s/kustomization.yaml} that lists them & \code{kubectl apply -k k8s/} \\
\bottomrule
\end{tabularx}
\end{center}

\begin{minted}{bash}
kubectl apply -f k8s/app.yaml     # one file
kubectl apply -k k8s/             # Kustomize
\end{minted}

Kustomize is part of \code{kubectl}, so staff need no other tool. With either
form, you write each object in full. This table shows an example Kustomize
layout. You can choose other file names.

\begin{center}
\small
\begin{tabularx}{\linewidth}{@{}lX@{}}
\toprule
File & Objects in the file \\
\midrule
\file{k8s/kustomization.yaml} & The list of files, and an optional \code{secretGenerator} \\
\file{k8s/postgres.yaml} & Postgres StatefulSet, headless Service, and the PVC template \\
\file{k8s/mcp.yaml} & MCP server Deployment and Service, with the probes \\
\file{k8s/config.yaml} & ConfigMap with \code{config.toml} \\
\file{k8s/inputs.yaml} & Input PVC for the target source \\
\file{k8s/networkpolicy.yaml} & NetworkPolicy for the agent Pod \\
\file{k8s/agent-test-pod.yaml} & Test Pod for the two test commands (Section~\ref{sec:testcmds}) \\
\file{k8s/run/agent-job.yaml} & Agent Job, with the init container, \code{securityContext}, and volumes \\
\bottomrule
\end{tabularx}
\end{center}

The agent Job is outside the main apply, because each run needs its own run ID
(step 6).

\textbf{Secrets.} Make the two Secrets in one of two ways. The first way is
\code{kubectl create secret generic} in the README, from environment variables
that staff set. The second way is a Kustomize \code{secretGenerator} that reads
a git-ignored \file{k8s/.env}. If you use the second way, the README tells staff
which variables to put in the file.

\begin{minted}{bash}
# Way 1: from environment variables that staff set
kubectl create secret generic llm-key --from-literal=LLM_API_KEY="$LLM_API_KEY"
kubectl create secret generic postgres --from-literal=POSTGRES_PASSWORD="$POSTGRES_PASSWORD"

# Way 2: k8s/kustomization.yaml reads k8s/.env, which git ignores
#   secretGenerator:
#     - name: llm-key
#       envs: [.env]
kubectl apply -k k8s/
\end{minted}

The secret names in this listing are examples. If staff cannot deploy your stack
from the README, your team loses the build points of
Section~\ref{sec:rubric}.

\section{Cluster setup}\label{sec:cluster}

\subsection{The staff cluster}\label{sec:kind}

All teams use the same cluster. The repository publishes it as
\file{cluster/kind-config.yaml}, with the pinned KIND version and the node
image. Use KIND v0.33.0. Test your stack on this cluster before you submit.

\begin{minted}{bash}
kind create cluster --config cluster/kind-config.yaml
\end{minted}

\filelisting{\file{cluster/kind-config.yaml}}{yaml}{../cluster/kind-config.yaml}

\subsection{CNI}\label{sec:cni}

% OPEN: feasibility task 3. The CNI choice and the Calico install command have no result on the build day.
The CNI is kindnet or Calico. Staff publish the choice by December 4. If the
choice is Calico, staff also publish the change to
\file{cluster/kind-config.yaml} and the install command.

Run a deny-all check before you trust your NetworkPolicy. Apply a policy that
denies all egress from one test Pod, then try a connection from that Pod. If the
connection succeeds, the CNI does not enforce NetworkPolicy.

\begin{minted}{bash}
kubectl run probe --image=busybox:1.37 --labels=app=probe -- sleep 3600
kubectl apply -f - <<'EOF'
apiVersion: networking.k8s.io/v1
kind: NetworkPolicy
metadata: {name: deny-all-egress}
spec:
  podSelector: {matchLabels: {app: probe}}
  policyTypes: [Egress]
EOF
kubectl exec probe -- nc -z -w 5 kubernetes.default.svc 443 && echo "NOT ENFORCED"
\end{minted}

The \code{nc} command must fail. A DNS failure is also a denial, because the
policy denies DNS too.

\subsection{Resources}\label{sec:resources}

% OPEN: feasibility task 8. The memory and CPU figures have no result on the build day.
The total memory and CPU for KIND, Postgres, the MCP server, and CLDK on a
laptop are to follow.

\subsection{Development path before December 4}\label{sec:devpath}

The lectures teach Jobs and Secrets on December 4. Until then, use this path:
\begin{itemize}
\item Run the agent as a plain Pod. Or run it on your laptop against Postgres
  through \code{kubectl port-forward}.
\item Keep the team key and the Postgres password in a \file{.env} file that git
  ignores.
\item Never put a key in a manifest, a ConfigMap, or a commit.
\end{itemize}

\subsection{Staff examples}\label{sec:examples}

Staff supply two short manifests with a comment on each field. They show the
pattern only. They are not the HW3 stack, and they do not use its names, images,
or ports.

\filelisting{\file{examples/k8s/statefulset.example.yaml}, a StatefulSet with a
headless Service}{yaml}{../examples/k8s/statefulset.example.yaml}

\filelisting{\file{examples/k8s/job-init.example.yaml}, a Job with an init
container}{yaml}{../examples/k8s/job-init.example.yaml}

\section{Deliverables}\label{sec:deliverables}
\begin{itemize}
\item All source code, a lockfile or equivalent that pins dependencies, and the
  commit of your HW2 baseline.
\item The Dockerfiles and all the Kubernetes manifests (Section~\ref{sec:staffsteps}).
\item \code{README.md}, as Section~\ref{sec:readme} describes, with the staff
  steps and the rubric index.
\item The unedited \code{report.json}, \code{trace.jsonl}, and ledger export for
  each run that your report cites.
\item The output of both test commands in the cluster.
\item The report (Section~\ref{sec:report}), as \code{REPORT.md} or a PDF.
\item A contribution statement that says who did what.
\item No API keys and no Secret values, in any file.
\end{itemize}

\section{Submissions}

We will create the HW3 assignment on CourseWorks (Canvas). Submit there and follow
the instructions on the assignment page. Each team submits once. Any member can
submit for the team.

\section{Academic integrity and AI tools}

You can use AI assistants to learn, debug your own code, review a design, or
complete small sections of code through autocomplete. Do not ask an agent to
complete the entire assignment and submit its output. You are responsible for
every line you submit. Every team member must be able to explain the entire
submission. The individual exams test this knowledge.

Do not use outside AI assistants to generate or change the target code. This
rule does not apply to the agents that you build for this assignment. They
change the workspace copy during the remediation stage. The course enforces its
AI-use policy strictly.

\end{document}
